\documentclass[14pt,a4paper]{extreport}
% preamble.tex — спільна преамбула для всіх лабораторних робіт. Підключається одразу після \documentclass у кожному lab*/main.tex.
%
% ЗБІРКА: LuaLaTeX (не pdfLaTeX). Причина: пакет listings під pdfLaTeX не читає UTF-8 і зупиняє збірку на кириличних коментарях у лістингах програмного коду. LuaLaTeX обробляє UTF-8 нативно.
%
% ШРИФТ: Liberation Serif — метрично сумісний з Times New Roman і містить повний набір кириличних гліфів (TeX Gyre Termes кирилиці не має, через що український текст не відображався взагалі).

\usepackage{fontspec}
\setmainfont{Liberation Serif}
\setmonofont{DejaVu Sans Mono}[Scale=0.82]

\usepackage{polyglossia}
\setmainlanguage{ukrainian}
\setotherlanguage{english}

\usepackage{microtype}

% Поля, інтервал, відступи й заголовки нижче взяті з шаблону НУЛП, за яким написана бакалаврська робота (../bachelor-thesis/thesis/main.tex). Той самий формат приймає В.М. Шиманський і для магістерських, тож усі звіти тримаємо в ньому, а не в тому, що переказували на консультації.
\usepackage[left=25mm,right=15mm,top=15mm,bottom=23mm,footskip=1.5mm]{geometry}
\usepackage{setspace}
% Саме \setstretch{1.5}, а не \onehalfspacing: останній дає близько 1.25.
\setstretch{1.5}

\usepackage{graphicx}
\usepackage{float}
\usepackage{booktabs}
\usepackage{array}
\usepackage{tabularx}
\usepackage{longtable}
\usepackage{enumitem}
% labelindent вирівнює маркер списку з абзацним відступом — так у шаблоні НУЛП.
\setlist{nosep}
\setlist[itemize]{label=\textbullet,leftmargin=*,labelindent=1.27cm}
\setlist[enumerate]{leftmargin=*,labelindent=1.27cm}
\usepackage{indentfirst}
\usepackage{amsmath}

\usepackage{hyperref}
\hypersetup{colorlinks=true,linkcolor=black,urlcolor=blue,citecolor=black}

% xurl дозволяє \url розривати адресу в будь-якому місці, а не лише після скісної риски. Без нього довгий сегмент шляху (напр. «scopus---podrobnaya- instruktsiya-dlya-uchenogo») виходить за межі сторінки, бо дефіс для hyperref не є точкою розриву. Вантажиться після hyperref — перевизначає його \UrlBreaks.
\usepackage{xurl}

% ---------------------------- Список використаної літератури ----------------------------
% biblatex + biber зі стилем IEEE: посилання «[1]» у тексті, записи — за порядком першої згадки. Обрано IEEE, а не gost-numeric (ДСТУ 8302:2015), — так вирішив автор звітів після порівняння обох варіантів.
% Кожен звіт із джерелами сам підключає свій references.bib через \addbibresource і друкує список \printbibliography[heading=none] під власним \section. Звіт без джерел від завантаженого biblatex нічого не втрачає.
% Оточення списку — стандартне для стилю: мітка від лівого поля, текст опису вирівняний після неї. Не підганяємо його під абзацний відступ нумерованих переліків тексту (labelindent у \setlist вище) — список літератури є окремим розділом зі своїм стандартним оформленням.
\usepackage[backend=biber,style=ieee]{biblatex}
% biblatex за замовчуванням додає розтяжний проміжок після «/ . : ? =» в адресах, і з вимкненими переносами рядок вирівнюється розрідженням самої адреси («https : / / www . scopus»). Рвати адресу й так дозволяє xurl, тож проміжок лише шкодить.
\setlength{\biburlbigskip}{0mu}
% Без додаткового інтервалу між описами: розділ має йти суцільним текстом, як решта звіту. Так само в шаблоні бакалаврської.
\setlength{\bibitemsep}{0pt}
\setlength{\bibparsep}{0pt}

% ---------------------------- Лістинги програмного коду ---------------------------- Використано minted (через Pygments), а не listings: listings не вміє коректно обробляти багатобайтові символи й «висмикує» кириличні коментарі з коду, зліплюючи їх в один рядок. Потребує прапорця -shell-escape.
\usepackage{xcolor}
\usepackage[newfloat]{minted}
% labelsep=period дає «Рис. 2.1. Назва» замість «Рис. 2: Назва» — формат НУЛП.
\usepackage[labelsep=period]{caption}

% baselinestretch=1 обов'язковий: без нього лістинг успадковує півторашку основного тексту, і довгий фрагмент коду перестає вміщатися у сторінку. Усередині [H]-флоата це не помилка збірки, а нескінченний цикл порожніх сторінок — lab1 так намолотив 161 тисячу сторінок.
\setminted{
  baselinestretch=1,
  fontsize=\footnotesize,
  linenos=true,
  frame=single,
  framesep=4pt,
  breaklines=true,
  breakanywhere=true,
  tabsize=4,
  numbersep=7pt,
  bgcolor=gray!8,
}
\setminted[python]{python3=true}

\SetupFloatingEnvironment{listing}{name=Лістинг,listname=Перелік лістингів}

% ---------------------------- Нумерація та оформлення заголовків ----------------------------
\usepackage{titlesec}

\setcounter{secnumdepth}{3}
\renewcommand{\thesection}{\arabic{section}}
\renewcommand{\thesubsection}{\arabic{section}.\arabic{subsection}}
\renewcommand{\thesubsubsection}{\arabic{section}.\arabic{subsection}.\arabic{subsubsection}}

% Розділ по центру, 16 пт, великими літерами; слово «РОЗДІЛ» опущене, бо звіт не дисертація. Номер лишається: \numberwithin прив'язує до нього нумерацію рисунків («Рис. 2.1»), і без нумерованих розділів вона б зламалася. \MakeUppercase четвертим аргументом отримує текст заголовка, тому в джерелі заголовки пишуться звичайним регістром.
\titleformat{\section}[block]
  {\centering\bfseries\fontsize{16}{19.2}\selectfont}
  {\thesection.}{0.3em}{\MakeUppercase}
% Підрозділи — зліва з абзацного відступу, звичайним регістром.
\titleformat{\subsection}[block]{\bfseries}{\thesubsection.}{1em}{}
\titleformat{\subsubsection}[block]{\bfseries\itshape}{\thesubsubsection.}{1em}{}

% Незіркові \titlespacing — обов'язкові. Без них titlesec застосовує зірковий варіант, який прибирає абзацний відступ у першого абзацу після заголовка й перекриває indentfirst. Значення — типові для report.
\titlespacing{\section}{0pt}{20pt}{6pt}
\titlespacing{\subsection}{1.27cm}{6pt}{0pt}
\titlespacing{\subsubsection}{1.27cm}{6pt}{0pt}

\setlength{\parindent}{1.27cm}
\setlength{\parskip}{0pt}
\renewcommand{\arraystretch}{1.3}

% Переноси слів вимкнені — так у шаблоні НУЛП. Натомість абзац розганяється міжслівними пробілами, тому tolerance і emergencystretch підняті: без них рядки виходили б за праве поле.
\hyphenpenalty=10000
\exhyphenpenalty=10000
\tolerance=9999
\setlength{\emergencystretch}{3em}

% ---------------------------- Підписи до рисунків і таблиць ----------------------------
\captionsetup[figure]{name=Рис.,justification=centering,font={bf,it}}
\captionsetup[table]{name=Таблиця,justification=raggedleft,singlelinecheck=false,labelfont=bf,textfont=normalfont}

% Рисунки, таблиці й формули нумеруються в межах розділу: «Рис. 2.1», а не «Рис. 5». Так у шаблоні НУЛП, і так само виглядали приклади на консультації.
\numberwithin{equation}{section}
\numberwithin{figure}{section}
\numberwithin{table}{section}

% ---------------------------- Колонтитули ----------------------------
\usepackage{fancyhdr}
\pagestyle{fancy}
\fancyhf{}
\fancyfoot[R]{\thepage}
\renewcommand{\headrulewidth}{0pt}
\fancypagestyle{plain}{%
  \fancyhf{}
  \fancyfoot[R]{\thepage}
  \renewcommand{\headrulewidth}{0pt}
}

% titlepage.tex — титульний аркуш, спільний для всіх дисциплін.
% Потребує визначених у subject.tex команд:
%   \SubjectName, \LecturerName, \StudentName, \StudentGroup, \AcademicYear
%
% Використання у тексті: \MakeLabTitle{1}{Тема роботи}
%
% Рід роботи. Дисципліни з практичними роботами перевизначають його після
% % preamble.tex — спільна преамбула для всіх лабораторних робіт. Підключається одразу після \documentclass у кожному lab*/main.tex.
%
% ЗБІРКА: LuaLaTeX (не pdfLaTeX). Причина: пакет listings під pdfLaTeX не читає UTF-8 і зупиняє збірку на кириличних коментарях у лістингах програмного коду. LuaLaTeX обробляє UTF-8 нативно.
%
% ШРИФТ: Liberation Serif — метрично сумісний з Times New Roman і містить повний набір кириличних гліфів (TeX Gyre Termes кирилиці не має, через що український текст не відображався взагалі).

\usepackage{fontspec}
\setmainfont{Liberation Serif}
\setmonofont{DejaVu Sans Mono}[Scale=0.82]

\usepackage{polyglossia}
\setmainlanguage{ukrainian}
\setotherlanguage{english}

\usepackage{microtype}

% Поля, інтервал, відступи й заголовки нижче взяті з шаблону НУЛП, за яким написана бакалаврська робота (../bachelor-thesis/thesis/main.tex). Той самий формат приймає В.М. Шиманський і для магістерських, тож усі звіти тримаємо в ньому, а не в тому, що переказували на консультації.
\usepackage[left=25mm,right=15mm,top=15mm,bottom=23mm,footskip=1.5mm]{geometry}
\usepackage{setspace}
% Саме \setstretch{1.5}, а не \onehalfspacing: останній дає близько 1.25.
\setstretch{1.5}

\usepackage{graphicx}
\usepackage{float}
\usepackage{booktabs}
\usepackage{array}
\usepackage{tabularx}
\usepackage{longtable}
\usepackage{enumitem}
% labelindent вирівнює маркер списку з абзацним відступом — так у шаблоні НУЛП.
\setlist{nosep}
\setlist[itemize]{label=\textbullet,leftmargin=*,labelindent=1.27cm}
\setlist[enumerate]{leftmargin=*,labelindent=1.27cm}
\usepackage{indentfirst}
\usepackage{amsmath}

\usepackage{hyperref}
\hypersetup{colorlinks=true,linkcolor=black,urlcolor=blue,citecolor=black}

% xurl дозволяє \url розривати адресу в будь-якому місці, а не лише після скісної риски. Без нього довгий сегмент шляху (напр. «scopus---podrobnaya- instruktsiya-dlya-uchenogo») виходить за межі сторінки, бо дефіс для hyperref не є точкою розриву. Вантажиться після hyperref — перевизначає його \UrlBreaks.
\usepackage{xurl}

% ---------------------------- Список використаної літератури ----------------------------
% biblatex + biber зі стилем IEEE: посилання «[1]» у тексті, записи — за порядком першої згадки. Обрано IEEE, а не gost-numeric (ДСТУ 8302:2015), — так вирішив автор звітів після порівняння обох варіантів.
% Кожен звіт із джерелами сам підключає свій references.bib через \addbibresource і друкує список \printbibliography[heading=none] під власним \section. Звіт без джерел від завантаженого biblatex нічого не втрачає.
% Оточення списку — стандартне для стилю: мітка від лівого поля, текст опису вирівняний після неї. Не підганяємо його під абзацний відступ нумерованих переліків тексту (labelindent у \setlist вище) — список літератури є окремим розділом зі своїм стандартним оформленням.
\usepackage[backend=biber,style=ieee]{biblatex}
% biblatex за замовчуванням додає розтяжний проміжок після «/ . : ? =» в адресах, і з вимкненими переносами рядок вирівнюється розрідженням самої адреси («https : / / www . scopus»). Рвати адресу й так дозволяє xurl, тож проміжок лише шкодить.
\setlength{\biburlbigskip}{0mu}
% Без додаткового інтервалу між описами: розділ має йти суцільним текстом, як решта звіту. Так само в шаблоні бакалаврської.
\setlength{\bibitemsep}{0pt}
\setlength{\bibparsep}{0pt}

% ---------------------------- Лістинги програмного коду ---------------------------- Використано minted (через Pygments), а не listings: listings не вміє коректно обробляти багатобайтові символи й «висмикує» кириличні коментарі з коду, зліплюючи їх в один рядок. Потребує прапорця -shell-escape.
\usepackage{xcolor}
\usepackage[newfloat]{minted}
% labelsep=period дає «Рис. 2.1. Назва» замість «Рис. 2: Назва» — формат НУЛП.
\usepackage[labelsep=period]{caption}

% baselinestretch=1 обов'язковий: без нього лістинг успадковує півторашку основного тексту, і довгий фрагмент коду перестає вміщатися у сторінку. Усередині [H]-флоата це не помилка збірки, а нескінченний цикл порожніх сторінок — lab1 так намолотив 161 тисячу сторінок.
\setminted{
  baselinestretch=1,
  fontsize=\footnotesize,
  linenos=true,
  frame=single,
  framesep=4pt,
  breaklines=true,
  breakanywhere=true,
  tabsize=4,
  numbersep=7pt,
  bgcolor=gray!8,
}
\setminted[python]{python3=true}

\SetupFloatingEnvironment{listing}{name=Лістинг,listname=Перелік лістингів}

% ---------------------------- Нумерація та оформлення заголовків ----------------------------
\usepackage{titlesec}

\setcounter{secnumdepth}{3}
\renewcommand{\thesection}{\arabic{section}}
\renewcommand{\thesubsection}{\arabic{section}.\arabic{subsection}}
\renewcommand{\thesubsubsection}{\arabic{section}.\arabic{subsection}.\arabic{subsubsection}}

% Розділ по центру, 16 пт, великими літерами; слово «РОЗДІЛ» опущене, бо звіт не дисертація. Номер лишається: \numberwithin прив'язує до нього нумерацію рисунків («Рис. 2.1»), і без нумерованих розділів вона б зламалася. \MakeUppercase четвертим аргументом отримує текст заголовка, тому в джерелі заголовки пишуться звичайним регістром.
\titleformat{\section}[block]
  {\centering\bfseries\fontsize{16}{19.2}\selectfont}
  {\thesection.}{0.3em}{\MakeUppercase}
% Підрозділи — зліва з абзацного відступу, звичайним регістром.
\titleformat{\subsection}[block]{\bfseries}{\thesubsection.}{1em}{}
\titleformat{\subsubsection}[block]{\bfseries\itshape}{\thesubsubsection.}{1em}{}

% Незіркові \titlespacing — обов'язкові. Без них titlesec застосовує зірковий варіант, який прибирає абзацний відступ у першого абзацу після заголовка й перекриває indentfirst. Значення — типові для report.
\titlespacing{\section}{0pt}{20pt}{6pt}
\titlespacing{\subsection}{1.27cm}{6pt}{0pt}
\titlespacing{\subsubsection}{1.27cm}{6pt}{0pt}

\setlength{\parindent}{1.27cm}
\setlength{\parskip}{0pt}
\renewcommand{\arraystretch}{1.3}

% Переноси слів вимкнені — так у шаблоні НУЛП. Натомість абзац розганяється міжслівними пробілами, тому tolerance і emergencystretch підняті: без них рядки виходили б за праве поле.
\hyphenpenalty=10000
\exhyphenpenalty=10000
\tolerance=9999
\setlength{\emergencystretch}{3em}

% ---------------------------- Підписи до рисунків і таблиць ----------------------------
\captionsetup[figure]{name=Рис.,justification=centering,font={bf,it}}
\captionsetup[table]{name=Таблиця,justification=raggedleft,singlelinecheck=false,labelfont=bf,textfont=normalfont}

% Рисунки, таблиці й формули нумеруються в межах розділу: «Рис. 2.1», а не «Рис. 5». Так у шаблоні НУЛП, і так само виглядали приклади на консультації.
\numberwithin{equation}{section}
\numberwithin{figure}{section}
\numberwithin{table}{section}

% ---------------------------- Колонтитули ----------------------------
\usepackage{fancyhdr}
\pagestyle{fancy}
\fancyhf{}
\fancyfoot[R]{\thepage}
\renewcommand{\headrulewidth}{0pt}
\fancypagestyle{plain}{%
  \fancyhf{}
  \fancyfoot[R]{\thepage}
  \renewcommand{\headrulewidth}{0pt}
}

\input{../../../../titlepage}
: \renewcommand{\WorkKind}{ПРАКТИЧНА РОБОТА}.
\newcommand{\WorkKind}{ЛАБОРАТОРНА РОБОТА}

\newcommand{\MakeLabTitle}[2]{%
\begin{titlepage}
% Титулка — не суцільний текст, а фіксована верстка з явних \\[...] проміжків,
% тож міжрядковий інтервал основного тексту (\setstretch{1.5} у preamble.tex)
% до неї не застосовний: він розтягував кожен із ~15 рядків на 20 %, і блок
% «Виконав» разом із «Львів --- рік» переповнював сторінку. Інтервал локальний
% для оточення titlepage, тому на решту документа це не впливає.
\setstretch{1}
\begin{center}
  {\large\bfseries МІНІСТЕРСТВО ОСВІТИ І НАУКИ УКРАЇНИ}\\[6pt]
  {\large\bfseries Національний університет «Львівська політехніка»}\\[6pt]
  {\large\bfseries Інститут комп'ютерних наук та інформаційних технологій}\\[6pt]
  {\large\bfseries Кафедра «Системи штучного інтелекту»}\\[1.2cm]

  \includegraphics[width=5cm]{../../../../logo}\\[1.2cm]

  {\Large\bfseries \WorkKind{} №#1}\\[0.5cm]
  з дисципліни\\[4pt]
  «\SubjectName»\\[0.6cm]
  на тему\\[4pt]
  \textit{#2}
\end{center}

% Відступ до блоку «Виконав» гнучкий, а не фіксований (було 2 cm). Фіксований
% разом із довгим родом роботи чи темою на два рядки переповнював сторінку, і
% «Львів --- рік» зі свого \vfill з'їжджав на окрему другу сторінку. \vfill
% стискається до нуля, коли місця бракує, тож переповнення стало неможливим.
\vfill

\begin{flushright}
  {\bfseries Виконав:}\\
  студент групи \StudentGroup\\
  \StudentName\\[1cm]
  {\bfseries Викладач:}\\
  \LecturerName\\
\end{flushright}

\vfill
\begin{center}
  Львів --- \AcademicYear
\end{center}
\end{titlepage}
\clearpage
}


% Метадані дисципліни. Використовуються титульним аркушем (titlepage.tex).

\newcommand{\SubjectName}{Науковий процес та професійна науково-дослідна робота}
% Два рядки: посада з науковим ступенем, далі прізвище. `\\` розриває рядок усередині flushright-блоку титульного аркуша. Посаду (професор, не доцент) підтвердив сам викладач.
\newcommand{\LecturerName}{проф. каф. СШІ д.т.н.,\\ Ізонін І.В.}
\newcommand{\StudentName}{Галич Максим}
\newcommand{\StudentGroup}{КНСШ-13}
\newcommand{\AcademicYear}{2026}

% Рід роботи для титульного аркуша — як у practice1 (див. коментар там).
\renewcommand{\WorkKind}{ЗВІТ\\[2pt]\normalfont\large про виконання практичної роботи}
% Рисунки — скриншоти, тобто рукотворний невідтворюваний вхід, тому лежать поряд із текстом і комітяться.
\graphicspath{{assets/}}
% Два файли джерел: власні джерела звіту та бібліотека Zotero. Другий перезаписує автоекспорт Better BibTeX, тому нічого рукотворного в ньому тримати не можна. Шляхи — відносно каталогу роботи, куди make render-report робить cd.
\addbibresource{references.bib}
\addbibresource{zotero.bib}

\usepackage{tcolorbox}
\tcbuselibrary{breakable}

% Один і той самий фрагмент тексту (demo-text.tex) верстається в кількох стилях цитування. biblatex зі спільної преамбули дає лише один стиль на документ, тому цитати й список у фрагменті формує рушій CSL напряму, через csl-demo.lua — з тих самих CSL-стилів і того самого zotero.bib, що й Zotero. Рамка відокремлює зразок від тексту звіту, а одинарний інтервал і менший кегль дають змогу вмістити його приблизно на сторінку.
\directlua{csldemo = dofile("csl-demo.lua")}
\newcommand{\StyleDemo}[3]{%
  \begin{tcolorbox}[breakable,colback=white,colframe=black!60,boxrule=0.5pt,arc=0pt]
    \small\setstretch{1}\setlength{\parindent}{1.27cm}%
    \directlua{csldemo.render("#1", "demo-text.tex", "zotero.bib", "Список літератури")}
  \end{tcolorbox}
  % \captionof поза оточенням caption-пакет не підтримує: без minipage він попереджає. hypcap вимкнено, бо поза float-ом немає його початку, куди ставити якір: посилання на рисунок веде до підпису.
  \noindent\begin{minipage}{\linewidth}\captionsetup{hypcap=false}\captionof{figure}{#2}\label{#3}\end{minipage}
}

% Рисунок-скриншот на всю ширину. Усі скриншоти роботи однакові за оформленням, тож оформлення в одному місці. Необов'язковий аргумент іде в \includegraphics — для обтинання (trim/clip) на етапі верстки, щоб оригінал скриншота лишався недоторканим.
\newcommand{\Screenshot}[4][]{%
  \begin{figure}[H]
    \centering
    \includegraphics[width=\linewidth,#1]{#2}
    \caption{#3}
    \label{#4}
  \end{figure}
}

\begin{document}

\MakeLabTitle{2}{Оформлення списку літератури з використанням систем управління бібліографічною інформацією}

\textbf{Мета роботи.} Встановити, налаштувати та навчитися використовувати безкоштовний бібліографічний менеджер Zotero для оформлення списку літератури у наукових роботах.

\section{Постановка завдання}

Інсталювати бібліографічний менеджер Zotero та навчитися ефективно його використовувати~\cite{guide}. Зокрема:

\begin{enumerate}
\item завантажити й встановити два варіанти бібліографічного менеджера Zotero~\cite{zotero} (десктопний і браузерний), зареєструватися та синхронізувати обидві версії;

\item записати обрану тему магістерської роботи, додати 2--4 абзаци тексту за її тематикою та поставити в них щонайменше шість посилань на релевантні темі наукові роботи, долучені до бібліотеки Zotero: по одній статті з журналу видавництва MDPI, видавництва Springer та видавництва Elsevier через розширення браузера; одну статтю з сервісу CEUR-WS.org вручну; по одному запису з баз PubMed та IEEE Xplore Digital Library за DOI;

\item сформувати бібліографію до цих робіт у стилі APA 7th edition;

\item змінити стиль на Vancouver та описати, як змінилися посилання в тексті та список літератури;

\item додати до бібліотеки новий стиль на власний вибір, замінити ним стиль Vancouver та описати зміни;

\item описати виконані кроки, проілюструвати їх рисунками, зробити висновок і описати проблеми, що виникли під час виконання роботи, та шляхи їх вирішення.
\end{enumerate}

Методичні рекомендації описують роботу із Zotero у текстовому процесорі MS Word. Звіти з дисципліни верстаються в LaTeX, тому текст із посиланнями та бібліографія тут формуються в LaTeX, а Zotero під'єднано до нього через розширення Better BibTeX (підрозділ~\ref{sec:latex}). Суть завдання — автоматичне формування посилань і списку літератури з бібліотеки Zotero та зміна стилю цитування — від цього не змінюється.

\section{Опис проведеної роботи}

\subsection{Встановлення та синхронізація Zotero}

Zotero існує у двох формах, які доповнюють одна одну: десктопна програма зберігає бібліотеку й формує бібліографію, а розширення браузера Zotero Connector розпізнає наукову публікацію на сторінці видавця і передає її метадані у програму. Десктопну програму~\cite{zotero} встановлено в Arch Linux з пакета \texttt{zotero-bin} репозиторію AUR, який містить офіційну збірку Zotero для Linux. Головне вікно програми наведено на рисунку~\ref{fig:zotero-desktop}: ліворуч — дерево бібліотеки з колекціями, у центрі — записи обраної колекції, а вгорі — вкладки з відкритими PDF-файлами статей, які Zotero дозволяє читати й анотувати без сторонньої програми.

% Знімок 3840x2336 px при 72 dpi, тобто 1 px = 1 bp. Обтинаємо порожні ділянки: знизу під списком записів (лишаємо верхні 1460 px) і праворуч порожню панель запису (лишаємо ліві 3126 px). Міняти при заміні скриншота.
\Screenshot[trim={0 876bp 714bp 0},clip]{zotero-desktop}{Головне вікно десктопної програми Zotero}{fig:zotero-desktop}

Розширення Zotero Connector версії 5.0.217 встановлено у браузер Zen, побудований на Firefox, з каталогу розширень Mozilla; на сторінці керування розширеннями воно увімкнене (рисунок~\ref{fig:zotero-connector}). Розширення працює лише разом із запущеною десктопною програмою: воно не має власного сховища і передає знайдені записи безпосередньо в бібліотеку програми.

\Screenshot{zotero-connector}{Розширення Zotero Connector у браузері}{fig:zotero-connector}

Для синхронізації зареєстровано обліковий запис на сайті Zotero з корпоративною поштовою адресою в домені \texttt{lpnu.ua} і виконано вхід у нього в налаштуваннях десктопної програми, у розділі \textit{Account} (рисунок~\ref{fig:zotero-sync}). Увімкнено автоматичну синхронізацію (\textit{Sync automatically}) разом із повнотекстовим вмістом документів, а прикріплені файли синхронізуються через сховище Zotero Storage. Після цього бібліотека зберігається не лише локально, а й на сервері Zotero, тож доступна з веббібліотеки та з інших комп'ютерів.

\Screenshot{zotero-sync}{Обліковий запис і налаштування синхронізації у десктопній програмі Zotero}{fig:zotero-sync}

\subsection{Формування бібліотеки}

Для джерел практичної роботи в бібліотеці створено окрему колекцію \textit{PR2}. Тема магістерської роботи — «Удосконалення методу синтезу мовлення з керуванням голосом за текстовим описом для української мови», тому всі шість джерел стосуються сучасного синтезу мовлення: клонування голосу, керування голосом за текстовою інструкцією та синтезу для мов з обмеженими даними. Джерела та способи їх додавання подано в таблиці~\ref{tab:sources}.

\begin{table}[H]
\caption{Джерела бібліотеки та способи їх додавання до Zotero}
\label{tab:sources}
\small
\begin{tabularx}{\linewidth}{|c|l|X|l|}
\hline
\textbf{№} & \textbf{Джерело} & \textbf{Публікація} & \textbf{Спосіб додавання} \\
\hline
1 & MDPI & Dorjee et al., Applied Sciences~\cite{dorjeeDiagnosingReferenceSideLanguage2026} & Zotero Connector \\
\hline
2 & Springer & Hanzlíček et al., LNCS~\cite{hanzlicekLeveragingZeroShotTTS2027} & Zotero Connector \\
\hline
3 & Elsevier & Kim, Choi, Computer Speech \& Language~\cite{kimParalinguisticAcousticFeature2027} & Zotero Connector \\
\hline
4 & CEUR-WS.org & Dzhaliuk et al., CPITS 2026~\cite{dzhaliukEvaluatingSpeakerVerification2026} & вручну \\
\hline
5 & PubMed & Bakkouche et al., Phonetica~\cite{bakkoucheWhatDeterminesSuccess2026} & за DOI \\
\hline
6 & IEEE Xplore & Lin et al., ICASSP 2026~\cite{linYouHearWhat2026} & за DOI \\
\hline
\end{tabularx}
\end{table}

\subsubsection{Додавання через розширення браузера}

Статті видавництв MDPI, Springer та Elsevier додано через Zotero Connector: на сторінці статті на сайті видавця з контекстного меню обрано пункт \textit{Zotero Connector} $\rightarrow$ \textit{Save to Zotero}, і Zotero зберіг запис у виділену в програмі колекцію \textit{PR2}. Меню пропонує кілька способів збереження тієї самої сторінки (рисунок~\ref{fig:connector-mdpi}), і вони відрізняються тим, звідки береться опис джерела:

\begin{itemize}
\item \textit{MDPI Journals} — транслятор, написаний саме під сайт видавця: він знає структуру сторінки й заповнює тип документа, авторів, назву, журнал, том, номер, сторінки та DOI;
\item \textit{Embedded Metadata} — метадані, які видавець вбудовує у заголовок сторінки для пошукових систем; зазвичай повні, але залежать від того, що саме видавець туди поклав;
\item \textit{DOI} — опис, отриманий з реєстру DOI за ідентифікатором, знайденим на сторінці;
\item \textit{Web Page with Snapshot} і \textit{Web Page without Snapshot} — збереження як звичайної вебсторінки, з копією сторінки або без неї; для наукової статті це неправильний тип, без авторів, журналу та DOI.
\end{itemize}

Обрано транслятор \textit{MDPI Journals}, бо він пропонується першим і дає найповніший опис. Збереження статті з журналу Applied Sciences видавництва MDPI наведено на рисунку~\ref{fig:connector-mdpi}.

\Screenshot{connector-mdpi}{Збереження статті видавництва MDPI~\cite{dorjeeDiagnosingReferenceSideLanguage2026} з контекстного меню Zotero Connector}{fig:connector-mdpi}

Так само збережено доповідь на конференції Text, Speech, and Dialogue, опубліковану видавництвом Springer у серії Lecture Notes in Artificial Intelligence (том 16940), підсерії Lecture Notes in Computer Science (рисунок~\ref{fig:connector-springer}), та статтю журналу Computer Speech \& Language видавництва Elsevier з платформи ScienceDirect (рисунок~\ref{fig:connector-elsevier}). Для сайту Springer меню пропонує той самий набір способів, лише з власним транслятором \textit{Springer Link}. Сайт позначає публікацію як \textit{Conference paper}, і транслятор зберіг її з відповідним типом запису Zotero \textit{Conference Paper}, а не \textit{Journal Article}, як статтю MDPI. Від типу залежить, які поля потрапляють в опис джерела: для доповіді це назва збірника праць, серія й видавець замість назви журналу, тому та сама стаття, збережена з неправильним типом, дала б у списку літератури хибний опис.

\Screenshot{connector-springer}{Збереження розділу збірника видавництва Springer~\cite{hanzlicekLeveragingZeroShotTTS2027} через Zotero Connector}{fig:connector-springer}

Статтю видавництва Elsevier збережено транслятором \textit{ScienceDirect}. Повний текст статті доступний лише за передплатою, проте це не завадило: транслятор бере опис зі сторінки статті, відкритої для всіх, тож запис у бібліотеці отримав повні метадані, хоча файл статті до нього не додано.

\Screenshot{connector-elsevier}{Збереження статті видавництва Elsevier~\cite{kimParalinguisticAcousticFeature2027} через Zotero Connector}{fig:connector-elsevier}

\subsubsection{Ручне внесення}

Статтю з праць семінару CPITS 2026, опублікованих у серії CEUR Workshop Proceedings~\cite{dzhaliukEvaluatingSpeakerVerification2026}, внесено вручну: у колекції створено новий запис типу \textit{Conference Paper} (доповідь на конференції) і заповнено його поля за титульною сторінкою статті та змістом тому. Заповнено назву доповіді, чотирьох авторів, чотирьох редакторів тому, назву праць, назву конференції, видавця, рік, місце проведення, номер тому, сторінки, серію, адресу файла статті та мову. Роль кожної особи задається окремо: редакторів внесено в той самий список, що й авторів, і змінено їм роль з \textit{Author} на \textit{Editor}. Картку запису із заповненими полями наведено на рисунку~\ref{fig:manual-ceur}.

Під час заповнення довелося розрізнити два схожі поля: \textit{Place} — це місце видавця, а \textit{Event Place} — місце проведення конференції. Семінар проходив у Києві, а видавець CEUR-WS.org місця видання не вказує, тому заповнено лише \textit{Event Place}, а \textit{Place} лишилося порожнім. Поле \textit{Citation Key} теж порожнє: ключ цитування присвоює розширення Better BibTeX (підрозділ~\ref{sec:latex}).

% Знімок вузький і високий (483x1629 px при 72 dpi, 1 px = 1 bp): на всю ширину він не вмістився б у сторінку, а [H]-рисунок, вищий за сторінку, заганяє LaTeX у нескінченний цикл порожніх сторінок. Тому ширина задана явно, а знизу обтято порожні поля після Language (лишаємо верхні 1592 px). Міняти при заміні скриншота.
\Screenshot[width=0.32\linewidth,trim={0 37bp 0 0},clip]{manual-ceur}{Запис статті з CEUR-WS.org~\cite{dzhaliukEvaluatingSpeakerVerification2026}, внесений вручну, із заповненими полями}{fig:manual-ceur}

Ручне внесення найтрудомісткіше і найбільше залежить від уважності: Zotero формує опис джерела лише з тих полів, які заповнено, тож пропущене поле так само зникає з бібліографії, а помилка в ньому переходить у неї без змін.

\subsubsection{Додавання за DOI}

Записи з баз PubMed та IEEE Xplore додано за DOI командою \textit{Add Item by Identifier} (піктограма чарівної палички на панелі інструментів). Для статті журналу Phonetica DOI взято з картки статті в PubMed (рисунок~\ref{fig:doi-pubmed}), для доповіді на конференції ICASSP~2026 — зі сторінки статті в IEEE Xplore (рисунок~\ref{fig:doi-ieee}). За DOI Zotero сам отримує повні метадані з реєстру DOI, тому цей спосіб найнадійніший: DOI однозначно ідентифікує публікацію, і ручне введення полів не потрібне. Крім DOI, те саме поле приймає ISBN книг, ідентифікатори PubMed (PMID), arXiv та ADS, тож статтю з PubMed можна було б додати і за її номером PMID~42287168.

\Screenshot{doi-pubmed}{Додавання статті з бази PubMed~\cite{bakkoucheWhatDeterminesSuccess2026} за DOI}{fig:doi-pubmed}

\Screenshot{doi-ieee}{Додавання статті з бази IEEE Xplore~\cite{linYouHearWhat2026} за DOI}{fig:doi-ieee}

Колекцію \textit{PR2} з усіма шістьма записами наведено на рисунку~\ref{fig:library}. Піктограма ліворуч від назви показує тип запису: три журнальні статті позначено аркушем, а три доповіді на конференціях — кафедрою. Праворуч від записів видно вкладення: до статей MDPI, IEEE та Phonetica Zotero додав PDF-файли, а до статті Elsevier, повний текст якої платний, — лише знімок сторінки.

Праворуч на рисунку відкрито картку доповіді Lin et al., доданої за DOI. Усі поля, зокрема назву праць, назву конференції, видавця, дату й сторінки, Zotero заповнив з реєстру DOI без жодного ручного введення. Водночас місце проведення конференції, Барселону, реєстр подав у полі \textit{Place}, тобто як місце видавця, тоді як для доповіді з CEUR-WS.org його свідомо внесено в \textit{Event Place}. Навіть автоматично отримані метадані варто переглядати: джерело даних може тлумачити поля інакше, ніж Zotero.

% Знімок 2000x1217 px при 72 dpi (1 px = 1 bp). Знизу обтято порожню частину списку разом із нижніми полями картки запису: лишаємо верхні 740 px, до поля Date включно. Міняти при заміні скриншота.
\Screenshot[trim={0 477bp 0 0},clip]{library-pr2}{Колекція \textit{PR2} бібліотеки Zotero з шістьма джерелами}{fig:library}

\subsection{Під'єднання Zotero до LaTeX}\label{sec:latex}

Для MS Word Zotero має власний плагін, а для LaTeX ту саму роль виконує розширення Better BibTeX~\cite{better-bibtex}. Воно присвоює кожному запису бібліотеки сталий ключ цитування, на який посилається текст, і експортує колекцію у файл формату BibLaTeX. Ключі утворено за типовою формулою розширення \texttt{auth.lower + shorttitle(3,3) + year}: прізвище першого автора, три перші значущі слова назви та рік, наприклад \texttt{linYouHearWhat2026}. Такий ключ довший за просте «прізвище й рік», зате однозначний навіть для кількох праць одного автора одного року. Колекцію \textit{PR2} експортовано у формат Better BibLaTeX — рідний для пакета biblatex, яким верстаються звіти: на відміну від BibTeX він має окремі поля для місця проведення конференції (\texttt{venue}), її назви (\texttt{eventtitle}) й повної дати (\texttt{date}). Експорт налаштовано з параметром \textit{Keep updated} у файл \texttt{zotero.bib} поряд із джерелом звіту (рисунок~\ref{fig:bbt-export}): після будь-якої зміни запису в Zotero файл оновлюється сам, тож бібліографія у звіті завжди відповідає бібліотеці. Нотатки й файли статей не експортуються, бо звіту вони не потрібні. Не ввімкнено й заміну назв журналів скороченнями: скорочену назву Better BibLaTeX і так експортує в окреме поле \texttt{shortjournal}, а стиль сам вирішує, яку з двох назв друкувати.

% Знімок діалогу малий (354x477 px), тому на всю ширину він вийшов би розмитим і завеликим.
\Screenshot[width=0.45\linewidth]{bbt-export}{Експорт колекції \textit{PR2} у формат Better BibLaTeX з автоматичним оновленням}{fig:bbt-export}

Стилі цитування в Zotero описано мовою Citation Style Language (CSL): кожен стиль — це XML-файл з правилами оформлення посилання в тексті та запису у списку літератури. Ті самі CSL-файли з офіційного репозиторію стилів~\cite{zotero-styles} вміє обробляти рушій citeproc-lua з пакета citation-style-language~\cite{csl-package}, який працює всередині LuaLaTeX. Тому текст за темою магістерської роботи записано в LaTeX з посиланнями на ключі бібліотеки, а цитати й список літератури до нього рушій формує автоматично, за обраним CSL-стилем. Нижче цей самий текст наведено в трьох стилях; змінюється лише назва стилю, а текст і бібліотека лишаються тими самими.

\subsection{Стиль APA 7th edition}

Стиль Американської психологічної асоціації сьомого видання~\cite{apa7} належить до стилів «автор--рік». Текст із посиланнями та сформований список літератури в цьому стилі наведено на рисунку~\ref{fig:apa}.

\StyleDemo{apa}{Текст із посиланнями та автоматично сформований список літератури у стилі APA 7th edition}{fig:apa}

Посилання в тексті має вигляд прізвища першого автора та року в круглих дужках: «(Lin et al., 2026)». Для праці двох авторів наводяться обидва прізвища через амперсанд — «(Kim \& Choi, 2027)», а для трьох і більше — лише перше зі скороченням «et al.». Список літератури не нумерується і впорядкований за алфавітом прізвищ перших авторів, тому порядок записів у ньому не збігається з порядком посилань у тексті. В описі джерела перелічено всіх авторів (до двадцяти) у формі «прізвище, ініціали», рік стоїть у дужках одразу після авторів, назва журналу та номер тому виділені курсивом, а DOI наведено як повну адресу \texttt{https://doi.org/\ldots}. APA вимагає писати назву статті з великої літери лише на початку, але стиль друкує назву так, як її збережено в Zotero: у трьох записах, які транслятори видавців MDPI, Springer та IEEE зберегли з великої літери в кожному слові, назви так і лишилися.

\subsection{Стиль Vancouver}

Стиль Vancouver склався в біомедичних журналах і нині підтримується Національною медичною бібліотекою США в посібнику Citing Medicine~\cite{citing-medicine}; саме його використовують статті, проіндексовані в PubMed. Після заміни стилю той самий текст і список літератури набули вигляду, наведеного на рисунку~\ref{fig:vancouver}.

\StyleDemo{vancouver}{Текст із посиланнями та автоматично сформований список літератури у стилі Vancouver}{fig:vancouver}

Зміни торкнулися і посилань у тексті, і списку літератури:

\begin{itemize}
\item посилання в тексті стали числовими: замість прізвищ і років — номери джерел у круглих дужках, від «(1)» до «(6)», тому речення тексту стали коротшими;
\item номери присвоєно за порядком першої згадки джерела в тексті, і список літератури тепер нумерований та впорядкований саме так, а не за алфавітом: стаття Lin et al., остання в алфавітному списку APA, стала записом 3, а стаття Dorjee et al., друга в ньому, — записом 5;
\item автори записані у формі «прізвище ініціали» без крапок і ком між ініціалами («Lin YC, Chou HC»), а з понад шести авторів наведено перших шість і «et al.» — так скорочено опис статті Bakkouche et al. з вісьмома авторами;
\item назви журналів більше не виділено курсивом; Vancouver друкує скорочену назву журналу, якщо її заповнено в полі \textit{Journal Abbr} запису, а в цих записах поле порожнє, тому лишилися повні назви («Applied Sciences»);
\item вихідні дані подано компактно: «2026 Jan [cited 2026 Oct 16];16(19):9466» — дата, том, номер у дужках і сторінка без слів-позначок;
\item діапазони сторінок скорочено до розрядів, що змінюються: «p.~278--89» замість «278--289», «p.~16472--6»;
\item для джерел з адресою в мережі додано позначку «[Internet]», дату звернення «[cited 2026 Oct 16]» з поля \textit{Accessed}, якщо її збережено (у внесеного вручну запису з CEUR-WS.org її немає), та адресу після «Available from:», а DOI стиль не наводить зовсім; доповідь Hanzlíček et al. адреси не має, тому описана як друкована.
\end{itemize}

Опис статті Bakkouche et al. не має тому й номера: після дати звернення одразу йде крапка з комою й адреса. Стаття опублікована в журналі Phonetica до виходу друкованого номера, тож тому й номера їй ще не присвоєно. Це не помилка стилю, а неповнота самих метаданих: після виходу номера запис у Zotero слід оновити за DOI.

\subsection{Додавання стилю IEEE Access}

Як новий стиль обрано IEEE Access — стиль однойменного журналу Інституту інженерів з електротехніки та електроніки (IEEE), відкритого мультидисциплінарного видання, в якому публікуються й роботи з обробки мовлення. Стиль знайдено в репозиторії Zotero Style Repository~\cite{zotero-styles} кнопкою \textit{Get Additional Styles} у розділі налаштувань \textit{Cite} і додано до переліку стилів (рисунок~\ref{fig:style-ieee}).

% Знімок 1704x1378 px при 72 dpi (1 px = 1 bp). Знизу обтято розділи Citation Dialog і Citation Options, які до стилів не стосуються: лишаємо верхні 1015 px, до рядка кнопок включно. Міняти при заміні скриншота.
\Screenshot[trim={0 363bp 0 0},clip]{style-ieee}{Стиль IEEE Access у переліку стилів цитування Zotero}{fig:style-ieee}

IEEE Access — залежний (\textit{dependent}) стиль. Його CSL-файл не містить жодного правила оформлення: лише назву журналу, ISSN і посилання на батьківський стиль IEEE~\cite{ieee-guide}, за яким журнал оформлює посилання. Так у репозиторії описано тисячі журналів, що користуються спільним стилем видавця: правила зберігаються один раз, у батьківському стилі, а залежні стилі лише вказують на нього. Тому, обравши IEEE Access, отримуємо оформлення IEEE; рушій CSL, як і Zotero, за посиланням у файлі сам підвантажує батьківський стиль. Текст і список літератури після заміни стилю Vancouver на IEEE Access наведено на рисунку~\ref{fig:ieee}.

\StyleDemo{ieee-access}{Текст із посиланнями та автоматично сформований список літератури у стилі IEEE Access}{fig:ieee}

IEEE Access, як і Vancouver, — числовий стиль, тому порядок джерел і їхні номери не змінилися. Змінилося інше:

\begin{itemize}
\item номер у тексті взято у квадратні дужки: «[1]» замість «(1)»;
\item у списку ініціали стоять перед прізвищем («Y.-C. Lin, H.-C. Chou»), останнього автора відокремлено сполучником «and», а для понад шести авторів наведено лише першого з «et al.» — «L. Bakkouche \textit{et al.}»;
\item назву статті взято в лапки, а назву журналу знову виділено курсивом; як і Vancouver, IEEE скорочує її лише за наявності поля \textit{Journal Abbr}, тому назви повні: «\textit{Applied Sciences}»;
\item вихідні дані розписано словами: «vol.~16, no.~19, p.~9466, Jan.~2026», а діапазони сторінок наведено повністю: «pp.~278--289»;
\item DOI повернувся до опису, але з префіксом «doi:», а не у формі адреси; адресу після «Available:» наведено лише для доповіді з CEUR-WS.org, яка DOI не має, а позначки «[Internet]» і дати звернення стиль не друкує.
\end{itemize}

\section{Результати}

У бібліотеці Zotero сформовано колекцію з шести наукових праць за темою магістерської роботи, додавши їх трьома способами: через розширення браузера, вручну та за DOI. Колекцію під'єднано до LaTeX через автоматичний експорт Better BibTeX, і текст за темою магістерської роботи з шістьма посиланнями зверстано в трьох стилях цитування. Відмінності стилів зведено в таблиці~\ref{tab:styles}.

\begin{table}[H]
\caption{Порівняння стилів цитування на прикладі сформованої бібліографії}
\label{tab:styles}
\small
\begin{tabularx}{\linewidth}{|l|X|X|X|}
\hline
\textbf{Ознака} & \textbf{APA 7th edition} & \textbf{Vancouver} & \textbf{IEEE Access} \\
\hline
Посилання в тексті & (Lin et al., 2026) & (3) & [3] \\
\hline
Порядок у списку & за алфавітом, без номерів & за першою згадкою, з номерами & за першою згадкою, з номерами \\
\hline
Автори & Lin, Y.-C., Chou, H.-C. & Lin YC, Chou HC & Y.-C. Lin, H.-C. Chou \\
\hline
Понад шість авторів & усі (до 20) & перші шість, et al. & перший, et al. \\
\hline
Назва журналу & повна, курсивом & без курсиву, скорочена за наявності \textit{Journal Abbr} & курсивом, скорочена за наявності \textit{Journal Abbr} \\
\hline
Вихідні дані & \textit{16}(19), 9466 & 2026 Jan [cited 2026 Oct 16];16(19):9466 & vol.~16, no.~19, p.~9466, Jan.~2026 \\
\hline
DOI & https://doi.org/\ldots & не наводиться, натомість адреса й дата звернення & doi: \ldots \\
\hline
\end{tabularx}
\end{table}

\section{Висновки}

У ході практичної роботи встановлено й синхронізовано десктопну та браузерну версії бібліографічного менеджера Zotero, сформовано бібліотеку з шести наукових праць за темою магістерської роботи та оформлено посилання на них і список літератури в стилях APA 7th edition, Vancouver та IEEE Access.

\begin{enumerate}[ref=\arabic*]
\item Опановано три способи додавання джерел, і вони суттєво різняться за надійністю. Додавання за DOI найнадійніше: метадані надходять з реєстру DOI і не залежать ні від сайту, ні від уважності користувача. Розширення браузера так само швидке, але покладається на транслятор конкретного сайту видавця. Ручне внесення найповільніше і єдине, в якому кожна помилка залежить від користувача, тому його варто лишати для джерел, які іншим способом не додаються.

\item Головна цінність менеджера — відокремлення даних про джерело від його оформлення. Опис кожної праці зберігається в бібліотеці один раз, а посилання в тексті та список літератури формуються з нього автоматично. Це прибирає ручну роботу з нумерацією, сортуванням і розділовими знаками, і саме в цій ручній роботі найчастіше виникають помилки оформлення.

\item Зміна стилю цитування зводиться до вибору іншого CSL-стилю: текст і бібліотека лишаються незмінними, а посилання й список перебудовуються повністю — від форми посилання в тексті до порядку записів. Для магістерської роботи це означає, що вимоги до оформлення списку літератури, які відрізняються в різних журналах і в самій кваліфікаційній роботі, можна виконати без переписування описів джерел.

\item Методичні рекомендації розраховані на MS Word, а звіти з дисципліни верстаються в LaTeX. Проблему вирішено розширенням Better BibTeX, яке присвоює записам сталі ключі цитування й автоматично оновлює файл бібліографії поряд із джерелом звіту, тож Zotero лишився єдиним місцем, де ведуться дані про джерела.

\item Спільна для всіх звітів преамбула використовує пакет biblatex, який дозволяє лише один стиль цитування на документ, а пакет citation-style-language разом із ним не підключається. Щоб показати той самий текст у трьох стилях в одному звіті, рушій CSL citeproc-lua викликано безпосередньо з LuaLaTeX: він обробляє ті самі CSL-стилі, що й Zotero.

\item У версії 0.9.1 рушія citeproc-lua виявлено помилку: для стилів, що визначають формат дати у своїй локалізації, як Vancouver, частини дати друкувалися без роздільника («2026June10»). Помилку виправлено перевизначенням відповідної функції рушія у модулі звіту, після чого дата має правильний вигляд «2026 June 10».

\item Якість бібліографії обмежена якістю метаданих, і навіть автоматично отримані метадані потребують перевірки. У сформованій бібліотеці виявлено чотири такі неточності: транслятори MDPI, Springer та IEEE зберегли назви статей з великої літери в кожному слові, і стиль APA, який вимагає великої літери лише на початку назви, друкує їх так; поле скороченої назви журналу (\textit{Journal Abbr}) лишилося порожнім, тому Vancouver та IEEE Access друкують повні назви журналів; транслятор MDPI зберіг дату статті Dorjee et al. як січень 2026~р., тоді як стаття вийшла у вересні; стаття в Phonetica, опублікована до виходу номера, не має тому й номера. Жодну з цих неточностей не виправити зміною стилю — лише в самому записі Zotero, звідки завдяки автоматичному експорту виправлення одразу потрапляє до звіту.
\end{enumerate}

\section{Список використаної літератури}

\printbibliography[heading=none]

\end{document}
